\documentclass[11pt]{article}
\usepackage{fontspec}
\setmainfont{Latin Modern Roman}
\usepackage{amsmath,amssymb,mathtools,unicode-math}
\setmathfont{Latin Modern Math}
\usepackage{xurl,hyperref}
\hypersetup{hidelinks,pdftitle={Problem statement}}
\newcommand{\sref}[2]{\hyperlink{source:#1:#2}{[S#2]}}
\newcommand{\cataloguescope}[1]{\paragraph{Mathematical question.}#1}
\begin{document}
\section{Open Problems in Integrable Systems — Argument shift method and Jordan-Kronecker invariants of Lie algebras}\label{statement}

\subsection{Definitions and mathematical statement}

\paragraph{Shared notation.}$\mathbb N=\{1,2,\ldots\}$, and $\mathbb Z,\mathbb Q,\mathbb R,\mathbb C$ denote the integers, rationals, reals and complex numbers. Any other conventions are specified locally.


For a finite-dimensional complex Lie algebra $\mathfrak g$ and $x,a\in\mathfrak g^*$, form the alternating pencil \[\mathcal A_{x+\lambda a}(u,v)=\langle x+\lambda a,[u,v]\rangle.\] Its Jordan--Kronecker decomposition is the simultaneous congruence canonical form of its pair of alternating forms: even-dimensional Jordan blocks at characteristic parameters and odd-dimensional Kronecker blocks. Its algebraic type records the number and sizes of Kronecker blocks and the sizes of Jordan blocks grouped by equal characteristic parameter; the numerical characteristic values themselves are ignored. For generic $(x,a)$ this type is the Jordan--Kronecker invariant of $\mathfrak g$.
\paragraph{Statement recorded in the supplied source.}
Characterize the algebraic types realizable by such generic argument-shift pencils as $\mathfrak g$ ranges over finite-dimensional complex Lie algebras. Determine all restrictions on the combinations of Jordan and Kronecker blocks and which admissible types can actually occur. \sref{statement}{2}
\subsection{Short English statement}
Which simultaneous alternating-form canonical types occur as generic argument-shift pencils of Lie algebras?
\subsection{Sources}
\begin{description}
\item[\hypertarget{source:statement:1}{S1}] ULAM.AI and the UnsolvedMath Contributors, UnsolvedMath: A Curated Collection of Open Mathematics Problems, catalogue entry. CC BY 4.0. \url{https://huggingface.co/datasets/ulamai/UnsolvedMath}.
\item[\hypertarget{source:statement:2}{S2}] A. Bolsinov, V. S. Matveev, E. Miranda and S. Tabachnikov, Open Problems, Questions, and Challenges in Finite-Dimensional Integrable Systems (2018), arXiv:1804.03737, Question 5.7 and generic Jordan-Kronecker definition. \url{https://arxiv.org/abs/1804.03737}.
\end{description}
\label{end:statement}
\end{document}
